\section{Research Plan and Roadmap}
\label{sec:plan}

Table~\ref{tab:roadmap} delineates planned research phases and risk mitigations addressing RQ3 and RQ4.

\begin{table}[t]
\caption{Structured Research Roadmap and Milestones}
\label{tab:roadmap}
\centering\scriptsize
\setlength{\tabcolsep}{2pt}
\begin{threeparttable}
\begin{tabular}{@{}lp{0.82in}p{1.02in}p{1.02in}@{}}
\toprule
Phase & Objective & Target Data & Analysis Method \\
\midrule
\textbf{P1} & Corpus \& Query Expansion (RQ3) & 100 repos (Java, TS, Python, Go); 50 JDs, RFPs, GitHub issues & Replicate $2 \times 2$ factorial design; mixed-effects ANOVA \\
\textbf{P2} & Graph-Enriched Headers (RQ3) & Call-graph \& import context vs. file-path headers & Paired Wilcoxon-Holm tests; ablation across header depths \\
\textbf{P3} & Human Verification Study (RQ4) & $N=20$ professional software developers & Task completion accuracy, inspection time, confidence \\
\bottomrule
\end{tabular}
\end{threeparttable}
\end{table}

\textbf{P1: Corpus \& Query Scaling (RQ3).} Expanding to 100 repositories across four languages (Java, TypeScript, Python, Go) and adding functional requirements from issues and RFPs alongside JDs. Language heterogeneity is mitigated by Tree-sitter parsers validated against syntax test suites.
\textbf{P2: Graph Context Enrichment (RQ3).} Prepending static call-graph and dependency context (signatures, imported interfaces) to test if relational graph context bridges the semantic gap for fine-grained AST chunks. Graph computation overhead is mitigated by precomputing indexes offline via SCIP/LSIF.
\textbf{P3: Human Practitioner Verification (RQ4).} Recruiting 20 developers to verify requirements with and without retrieved code evidence, measuring task accuracy, inspection time, and confidence. Subject variance is mitigated via a counterbalanced Latin-square design with skill stratification.
